  % !TeX root = book.tex
\chapter{What's in it for other communities?}\label{sec:for-other-communities}

\vspace{1cm}

The long-term vitality and reach of the Prolog Trinity ecosystem depends on its ability to attract and serve a broader constellation of users beyond the traditional Prolog community. Within this ecosystem, logic programming researchers gain a venue in which their ideas can not only be demonstrated in a suitable GUI, but also \textit{used} in practice through remote calls from external applications. Researchers and practitioners in the Semantic Web and Big Data worlds may find in Web Prolog a natural complement to existing technologies, since Prolog is well suited to data preprocessing, interactive construction of small, composable views, and transparent integration of heterogeneous sources such as relational databases. Symbolic and neuro-symbolic AI developers can exploit the Prolog Web's declarative expressiveness, agent-based communication model, and support for distributed inference, using Prolog both as a foundational language for symbolic AI and as a common representation for rules extracted from sub-symbolic predictors or for constraints imposed on machine-learning models to mitigate bias and opacity. Educators, system architects, and developers of intelligent web applications may likewise be drawn to an ecosystem that excels at building expert systems and logic solvers (probabilistic, abductive, or inductive) over heterogeneous data, particularly as it matures and integrates more tightly with established Web protocols and standards. In this way, dogfooding by the core community not only strengthens the platform from within but also demonstrates its readiness to welcome and support a diverse range of users.

Figure~\ref{fig:prolog-web-intro-20} introduces symbols that hint at neighbouring ecosystems. The emblem with nested red and black circles at node \texttt{n4.org} is the logo for \textit{cplint}, a probabilistic logic programming platform. The blue database cylinders represent RDF triple stores, standing in for Semantic Web back-ends that can be exposed as Prolog Web nodes. On the right, the familiar Python logo appears together with the OpenAI emblem, indicating a Python-based service that fronts a large language model such as ChatGPT and makes it available to other agents via a FFI. 

%The diagram in Figure~\ref{fig:prolog-web-intro-20} adds new symbols to Figure~\ref{fig:prolog-web-intro-200} that suggests some of these broader use cases.

%The long-term vitality and reach of the Prolog Trinity ecosystem depends on its ability to attract and serve a broader constellation of users. Among these are researchers and practitioners in the Semantic Web community, who may find in Web Prolog a natural complement or alternative to existing logic-based semantic technologies. Likewise, developers working in symbolic and neuro-symbolic AI stand to benefit from the Prolog Web's declarative expressiveness, agent-based communication model, and support for distributed inference. Educators, system architects, and developers of intelligent web applications may also be drawn to the ecosystem, particularly as it matures and integrates with established web protocols and standards. In this way, dogfooding by the core community not only strengthens the platform from within but also signals its readiness to welcome and support a diverse range of users beyond the Prolog world.
%
%
%In the realms of Big Data, the Semantic Web, and Artificial Intelligence, 
%Prolog offers significant strengths in the domains of Big Data and AI. For Big Data, Prolog, especially when enhanced with tabling support, is well-suited for data preprocessing and can transparently access and integrate information from diverse sources, such as relational database management systems. It enables the creation of data views through small, composable predicates that can be debugged interactively and separately. Systems like SWISH datalab leverage Prolog for cooperative development of data preprocessing and statistical analysis, while other Prolog implementations like Yap are designed to handle large datasets effectively. In AI, Prolog is a foundational concept in symbolic AI, offering solutions for critical issues in data-driven AI, such as inherent biases and opacity. It can serve as a common language for symbolic rules extracted from sub-symbolic predictors and can impose constraints on machine and deep learning models to reduce biases. The integration of sub-symbolic AI can also enhance Prolog's mechanisms, as seen in the use of neural networks to improve the efficiency of inductive logic programming (ILP) algorithms. Neuro-symbolic approaches further combine Prolog's inferential capabilities with neural networks' pattern-matching abilities to create more intelligent systems. Furthermore, Prolog is ideal for developing novel expert systems and logic solvers that utilize probabilistic, abductive, or inductive inference for symbolic AI applications, and it excels at integrating and reasoning over heterogeneous data.

\todo{TODO: Caption must be fixed}

\begin{figure}[h]
    \centering
	\includegraphics[width=\textwidth]{prolog-web-intro-20.png}
    \caption{The Prolog Trinity: End-user clients in interaction with Prolog agents on the Prolog Web. (The logotype with nested red and black circles is a system for probabilistic logic programming.)}
    \label{fig:prolog-web-intro-20}
\end{figure}

\noindent These additions suggest that Prolog nodes are intended to live in the same networked space as Semantic Web infrastructure and neural/ML services. Section~\ref{sec:broader-LP-community}, \ref{sec:semantic-web-community} and~\ref{sec:ai-community} return to these connections from the perspectives of the Logic Programming, the Semantic Web and AI communities respectively. What front-end programmers may gain is discussed in Section~\ref{sec:for-web-programmers}. Last but not least, although Erlang is not explicity represented in Figure~\ref{fig:prolog-web-intro-20}, it is covered in Section~\ref{sec:erlang-community}.

\section{Toward a Logic Programming Web}
\label{sec:lp-web-vision}

Earlier chapters have treated the Prolog Web primarily as a response to the current state of Prolog and its ecosystem, and as a concrete way of putting the Prolog Trinity into practice. In this section we take a broader view and ask what the proposal might offer to the wider logic programming community, including traditions that do not identify themselves as ``Prolog systems'' in any narrow sense. If you work on Answer Set Programming, Datalog variants, constraint systems, probabilistic logic programming, or other LP formalisms, what is in it for you?

The short answer is: a Web-native integration layer that your system can join with minimal effort. Section~\ref{sec:lp-landscape} surveyed the diversity of the logic programming landscape; Section~\ref{sec:one-prolog} argued that the Prolog Web's profile hierarchy can accommodate that diversity without forcing any system to adopt features foreign to its paradigm. Here we demonstrate this concretely.

As illustrated in Figure~\ref{fig:10.1}, several research systems can already be deployed on the Prolog Web. The integration paths differ in detail but share a common pattern: wrap the system's query interface in the Web Prolog protocol, advertise the appropriate profile, and participate.

Consider probabilistic logic programming. The cplint system, introduced in Section~\ref{sec:lp-landscape}, is implemented as a SWI-Prolog library and is already accessible through ``cplint on SWISH.'' Because SWISH builds on Pengines, cplint is in practice already reachable through the Pengines protocol -- and therefore through the subset of Web Prolog that Pengines implements. The following example illustrates this directly.

\subsection*{Pengines calling cplint}

Suppose you have two urns: urn1 contains 40 blue balls and 20 red balls and urn2 contains 25 blue balls and 30 red balls. You throw an unbiased coin and, if it turns out head, you draw a ball from the first urn, if it turns out tails you draw a ball from the second urn. Write a program modeling this process and a query for answering the question ``What is the probability of drawing a blue ball?''

\begin{verbatim}
?- pengine_rpc('http://cplint.ml.unife.it', prob(blue,Prob), [
     src_text("
       :- use_module(library(pita)).
       :- pita.
       :- begin_lpad.
       blue : 2/3; red : 1/3 :- urn1.
       blue : 25/55; red : 30/55 :- urn2.
       urn1 : 0.5 ; urn2 : 0.5.
       :- end_lpad.
     ")
   ]).
Prob = 0.5606060606060606.
?-
\end{verbatim}


\noindent
LPS, also implemented in SWI-Prolog and SWISH, can be accessed in the same manner. However, as it now stands, the result we receive is a large term (a dict) that serves as instructions for how to display the result graphically.

ProbLog presents a different integration path. As discussed in Section~\ref{sec:lp-landscape}, ProbLog~2 is a Python-based toolbox with a strictly declarative execution model: it does not support side-effecting operations such as I/O or dynamic modification of the clause database. Rather than being a limitation, this makes ProbLog an ideal candidate for an ISOBASE node, since ISOBASE assumes nothing more than a pure ISO-style toplevel without mutable state or Erlang-style concurrency. Wrapping the ProbLog runtime as an ISOBASE node would allow other agents to submit goals, obtain answers, and incorporate probabilistic reasoning results into larger workflows on the Prolog Web.

Multi-paradigm languages such as Picat offer yet another integration path. Picat's constraint solving, planning, and optimisation capabilities could be exposed as specialised services on the network, addressable by any client that can issue a Web Prolog query.

The common pattern is that participation does not require rewriting a system from the ground up. It requires a thin protocol layer -- a wrapper that translates between the Web Prolog message format and the system's native query interface -- and a declaration of the appropriate profile. Systems built on SWI-Prolog are closest to immediate deployment; systems with their own runtimes need a lightweight adapter, but the protocol is deliberately simple enough to make this practical.

The broader point is architectural. The logic programming community has produced a remarkable diversity of formalisms and implementations, but they remain largely isolated from one another. A researcher using ProbLog cannot easily delegate a constraint satisfaction sub-problem to a Picat node, or combine probabilistic inference with ASP-style model generation, without building bespoke integration code. The Prolog Web offers a common substrate -- a shared protocol, a shared addressing scheme, and a shared notion of ``ask a question, get answers'' -- that could turn isolated systems into a federated network. Whether such a Logic Programming Web will emerge depends on adoption by these communities, but the infrastructure described in this book is designed to make that adoption as frictionless as possible.


\section{For the broader logic programming community? -- OLD}\label{sec:broader-LP-community}

Earlier chapters have treated the Prolog Web primarily as a response to the current state of Prolog and its ecosystem, and as a concrete way of putting the Prolog Trinity into practice. In this section we take a broader view and ask what the proposal might offer to the wider logic programming (LP) community, including traditions that do not identify themselves as ``Prolog systems'' in any narrow sense. If you work on Answer Set Programming, Datalog variants, constraint systems, probabilistic logic programming, or other LP formalisms, what is in it for you?

The architectural vision behind Web Prolog and the Prolog Web is not limited to the Prolog tradition. It points toward the possibility of a Logic Programming Web in a broader sense. 

Most logic programming systems take advantage of the unique Prolog features. LPS, already discussed in Chapter 7, certainly does.

As illustrated in Figure~\ref{fig:prolog-web-intro-20}, research areas such as probabilistic logic programming can already be deployed on the Prolog Web: systems such as \textit{cplint}, \textit{ProbLog}, and multiparadigm languages such as \textit{Picat} can be wrapped in a Web Prolog--compatible interface and immediately participate. These examples demonstrate that the Prolog Web need not be confined to classical Prolog, but can serve as an integrating substrate for a wide range of logic programming technologies.

\begin{itemize}
\item \textbf{cplint} is a suite of tools for probabilistic logic programming, implemented as a Prolog library and available in particular for SWI-Prolog. It supports both inference and learning for Logic Programs with Annotated Disjunctions (LPADs) and related formalisms under the distribution semantics, allowing users to define discrete probability distributions and continuous probability densities directly in Prolog. cplint is accessible both as a local Prolog pack and through ``cplint on SWISH,'' a web-based environment where probabilistic logic programs can be edited, queried, and learned from data via a browser.
\item \textbf{ProbLog} is a probabilistic logic programming language that extends Prolog with probabilistic facts, allowing each fact to hold with a given probability and thereby combining logical atoms with random variables under the distribution semantics. A ProbLog program defines a probability distribution over possible worlds, and queries return the probability that a given atom is true rather than just a yes/no answer. The current ProbLog 2 system is a Python-based toolbox offering efficient algorithms for probabilistic inference, most probable explanations, sampling, and parameter learning, and is available both as a library and via online tools from the DTAI group at KU Leuven. 
\item \textbf{Picat} is a multi-paradigm, Prolog-like logic-based programming language that integrates logic programming, constraint solving, dynamic programming, planning, and scripting in a single framework. It offers several specialized solver modules and a high-level modeling style where problems are expressed declaratively and solved via built-in search and optimization mechanisms. Picat's pattern-matching rules, tabling, and constraint-based primitives make it well suited for combinatorial search, AI planning, and optimization tasks, while its familiar syntax and tight integration of different paradigms aim to provide a practical alternative to both traditional Prolog systems and dedicated constraint or planning languages.
\end{itemize}

\subsection*{Pengines calling cplint}

\begin{quotation}
Suppose you have two urns: urn1 contains 40 blue balls and 20 red balls and urn2 contains 25 blue balls and 30 red balls. You throw an unbiased coin and, if it turns out head, you draw a ball from the first urn, if it turns out tails you draw a ball from the second urn. Write a program modeling this process and a query for answering the question ``What is the probability of drawing a blue ball?''
\end{quotation}

\begin{verbatim}
?- pengine_rpc('http://cplint.ml.unife.it', prob(blue,Prob), [
       src_text("
           :- use_module(library(pita)).
           :- pita.
        
           :- begin_lpad.
        
           blue : 2/3; red : 1/3 :- urn1.
           blue : 25/55; red : 30/55 :- urn2.
        
           urn1 : 0.5 ; urn2 : 0.5.
        
           :- end_lpad.
       ")
   ]).
Prob = 0.5606060606060606.
?-
\end{verbatim}

\noindent LPS, also implemented in SWI-Prolog and SWISH, can be accessed in the same manner. However, as it now stands, the result we receive is a big term (an dict) that serve as instructions for how to display the result graphically.

A key enabler of this integration is the node profile architecture introduced in earlier chapters. Since not all systems support the same operational features, the profiles ensure that each system participates at a level consistent with its expressive capabilities. This becomes particularly clear in the case of \textit{ProbLog}. Designed around the distribution semantics, ProbLog maintains a strictly declarative execution model and does not support side-effecting operations such as I/O or dynamic modification of the clause database. Rather than being a limitation, this makes ProbLog an ideal candidate for an \textsc{ISOBASE} node, since \textsc{ISOBASE} assumes nothing more than a pure ISO-style toplevel without mutable state or Erlang-style concurrency. Wrapping the ProbLog runtime as an \textsc{ISOBASE} node allows other agents to submit goals, obtain answers, and incorporate probabilistic reasoning results into larger workflows on the Prolog Web. More feature-rich systems may adopt the \textsc{ISOTOPE} or \textsc{ACTOR} profiles, but languages such as ProbLog fit naturally and cleanly into the \textsc{ISOBASE} profile.

A major design motive is precisely this: to enable existing Prolog systems -- diverse in capabilities, internal representations, and implementation philosophies -- to communicate with one another using Web Prolog as a \textit{lingua franca}. By presenting each system as a node that exposes a uniform Web Prolog interface, heterogeneous engines can exchange data, script one another's toplevels and actors, and cooperate in distributed computations. Extending this principle, any system for which such an interface can be provided -- deductive databases, non-Prolog logic programming languages, or systems written in Erlang, Java, JavaScript, or Python -- can join the same conversational fabric by adopting the profile appropriate to its operational model.

From the perspective of the broader logic programming community, the Prolog Trinity ecosystem thus offers a concrete path toward greater interoperability and visibility. Over the years LP has developed a rich variety of languages and reasoning paradigms, including Answer Set Programming, Datalog variants, constraint systems, and probabilistic frameworks. Yet these systems often remain isolated, each with its own tools, formats, and runtimes. A logic-enabled Web of communicating agents provides a unifying context in which these approaches can interact without sacrificing their distinct strengths. One might imagine Web Prolog agents consulting ASP solvers for combinatorial tasks, invoking Datalog engines over large datasets, or delegating probabilistic inference to specialised systems such as ProbLog or \textit{cplint} -- each wrapped as a Prolog Web node under the profile suitable to its capabilities.

This book does not attempt to contribute to the logical foundations of logic programming itself; instead it focuses on extra-logical infrastructure: the practical mechanisms that allow diverse systems to coexist, interoperate, and be deployed on the Web. In that sense, the Trinity is not merely a proposal for the Prolog community narrowly conceived. It is an invitation to the entire logic programming field to anchor its research prototypes, experimental systems, and mature engines in a shared, web-native environment, where ideas can circulate more freely, capabilities can be reused across implementations, and logic programming can present a more coherent face to the wider world.

%Figure~14.1 also introduces symbols that hint at neighbouring ecosystems. The logo with nested red and black circles at node \texttt{n4.org} denotes \textit{cplint}, a probabilistic logic programming platform. The blue database cylinders with small linked nodes represent RDF triple stores, standing in for Semantic Web back-ends that can be exposed as Prolog Web nodes. On the right, the familiar Python logo appears together with the green OpenAI emblem, indicating a Python-based service that fronts a large language model such as ChatGPT and makes it available to other agents via a web API. These additions suggest that Prolog nodes are intended to live in the same networked space as Semantic Web infrastructure and neural/ML services; Sections~14.2 and~14.3 return to these connections from the perspectives of the Semantic Web and AI communities respectively

%Potentially even a Logic Programming Web more broadly might be possible. For example, as suggested in Figure~\ref{fig:prolog-web-intro-20}, some of the more recent research in the area of logic programming can be deployed on the Prolog Web, such as research into probabilistic reasoning. The logotype with nested red and black circles is a probabilistic logic programming language known as \textit{cplint}.\footnote{https://cplint.ml.unife.it} Other systems that might be easily deployed on the Prolog Web are \textit{Problog},\footnote{\url{https://dtai.cs.kuleuven.be/problog/}} and perhaps a language such as \textit{Picat}.\footnote{\url{https://www.picat-lang.org}}
%
%---
%
%\noindent In Figure~\ref{fig:prolog-web-intro-20}, \includegraphics[width=0.25cm]{cplint-logo.png}
%is the logo for a probabilistic logic programming language known as \textit{cplint}. This is just one example of a logic programming platform that is not Prolog. There are many others.
%
%In the figure, \includegraphics[width=0.25cm]{rdf-logo.png} stands for a triple store.
%
%\includegraphics[width=0.25cm]{python-logo.png}
%is the logo for the Python programming language and  The logo \includegraphics[width=0.25cm]{openai-logo.png} stands for OpenAI.
%
%---
%
%\noindent A major motive behind the proposed design of the Web Prolog language and the architecture of the Prolog Web is to allow the many Prolog systems currently in existence to talk to each other using Web Prolog as a \textit{lingua franca}. As an extension of this idea we also propose that, as long as we can wrap an interface consistent with Web Prolog around them, such conversations might involve not only Prolog systems, but also many non-Prolog systems, such as (deductive) databases, other logic programming systems, as well as systems running other programming languages such as Erlang, Java or JavaScript. 
%
%Using Web Prolog as an embedded scripting language would allow these systems to exchange information using Web Prolog as a \textit{lingua franca}, allow clients to script toplevels and other actors running on them, and allow us to integrate all these systems into a truly programmable Prolog Web.
%
%---
%
%The Prolog Trinity is not merely an evolution for Prolog; it is an invitation to the entire Logic Programming (LP) community to embrace a new era of relevance and impact. For decades, the LP paradigm has offered profound insights into computation, knowledge representation, and reasoning. From foundational theoretical work to the development of diverse LP languages (beyond Prolog, including Answer Set Programming, Datalog, and others), the community has cultivated a rich intellectual heritage. However, the translation of these potent ideas into mainstream software development and widely adopted applications has often faced headwinds.
%
%---
%
%While Web Prolog is the initial focus, the architectural vision of a Prolog Web -- a web of intercommunicating logical agents and services -- is expansive enough to potentially embrace and interoperate with other forms of logic programming. Imagine Web Prolog agents that can query Answer Set Programming solvers for complex combinatorial problems or interact with Datalog engines for efficient data-intensive reasoning. The Trinity can serve as an integrating fabric, allowing the diverse strengths of the entire LP ecosystem to be brought to bear.
%
%---
%
%
%\noindent This book does not purport to contribute to the logic programming aspects of Prolog, only to its extra-logical aspects. Still, the Prolog Web could presumably work as a valuable service to the logic programming community at large, allowing members to present their work through APIs that expose their research systems to the Prolog Web for others to easily try out and perhaps also use, regardless of if it is written in Prolog or not. Indeed, it might inspire logic programming researchers to develop proof-of-concept implementations that work well with the Prolog Web, and, after some time gathering real life experience, maybe some of the things that they do can be further standardized as new profiles of Web Prolog. In the best of worlds, this might cause the Prolog Web to evolve organically.
%
%
%---
%
%
%Note that a ProbLog node only needs ISOBASE or ISOTOPE capabilities.
%
%Web Prolog thus becomes a language for connecting computational logic systems of many different kinds - deductive databases, probabilistic systems, abductive systems, always with terms that are ISO.
%
%Such systems are not necessarily written in Prolog
%
%Neural networks
%
%compare with micro services
%
%orchestrated by programs written in Web Prolog perhaps in combination with SCXML v 2.0

\section{For the semantic web community?}\label{sec:semantic-web-community}

Chapter~5 argued that Web Prolog and the Prolog Web can be positioned as a kind of \emph{semantic-web logic programming language} and a complementary ``third tower'' alongside RDF, OWL, and SPARQL. Rather than rehearse those arguments, this section looks at the same picture from the point of view of researchers and practitioners already committed to the Semantic Web stack. The question is simple: if you work with RDF graphs, OWL ontologies, and SPARQL endpoints, what is in it for you?

A first answer is that the Prolog Trinity ecosystem offers a missing \emph{programming layer} for the Semantic Web -- a layer where logic and action are integrated, and where semantic applications can not only query and infer, but also react, maintain state, and communicate. RDF, OWL, and SPARQL provide a powerful foundation for representing and querying knowledge, but they stop short of providing a uniform, web-native language in which semantic agents can be written. Web Prolog fills precisely that gap. It gives Semantic Web developers a logic-based scripting language in which they can express rules, policies, and workflows that live \emph{on top of} RDF data and ontologies, instead of having to embed such logic in ad-hoc JavaScript, Java, or Python code around a triple store.

A second answer is that the Trinity supplies a concrete and deployable \emph{agent model} for the Semantic Web. Chapter~7 introduced the Agentic Web as a vision of a Web populated not only by linked datasets but by software agents capable of perceiving, reasoning, and acting. The Prolog Web realises a fragment of that vision: Web Prolog agents can consume RDF data, consult OWL ontologies, perform symbolic inference, and then publish new knowledge or trigger follow-on actions. For Semantic Web practitioners who have long spoken about ``semantic agents'' in rather abstract terms, the Trinity offers an explicit architecture and a programmable substrate for building them.

Third, the Prolog Web gives the Semantic Web community a practical way to connect heterogeneous reasoning systems without abandoning the existing stack. Prolog nodes may be backed by persistent RDF triple stores, relational databases, graph databases, or specialised reasoners such as probabilistic or abductive engines. From the outside, these appear as Prolog Web nodes with appropriate profiles, exposing logical APIs that exchange ISO-compatible terms.

Seen from this perspective, Web Prolog is not a rival to RDF, OWL, or SPARQL, but a way of harvesting decades of logic-programming research to make the Semantic Web more usable. Features such as tabling, constraint handling rules, and well-founded semantics already exist in mature Prolog systems; the Prolog Trinity ecosystem proposes to reuse this infrastructure rather than re-implement it in bespoke rule engines attached to triple stores. For Semantic Web practitioners, this means that advanced reasoning facilities -- non-monotonic reasoning, or domain-specific rule languages -- can be prototyped and deployed on top of existing data.

The architectural side of the story is equally important. The Prolog Web provides a disciplined way of deploying such logic-based components as nodes in a distributed system that speaks standard Web protocols. Nodes can front-end SPARQL endpoints, serve or consume RDF, or participate in Linked Data workflows, while also hosting agents that communicate asynchronously and carry out higher-level tasks. This suggests a path toward what one might call a \emph{Semantic Prolog Web}: a hybrid in which RDF and OWL continue to carry the semantic load, while Web Prolog and Prolog agents provide the programmable glue, orchestration, and long-lived behaviour that many applications need.

Finally, there is an institutional and standardisation angle. The W3C Semantic Web activity has produced an impressive stack of standards for data and ontology representation, but has so far remained largely programming-language neutral. The Trinity proposal invites the Semantic Web community to consider whether it might be time to recognise a web-native logic programming language, developed in concert with the ISO Prolog community, as a first-class tool in its ecosystem. Supporting the development and standardisation of Web Prolog, the Prolog Web, and Prolog agents would not replace the existing stack; it would supply a coherent way of turning its machine-understandable semantics into running, distributed, semantically aware applications.

If that invitation is accepted, the benefits are mutual. The Semantic Web gains a practical, logically grounded programming and agent layer; the Prolog Trinity gains a natural home in an established standards landscape. Together they move the Web one step closer to a world in which linked data is not only published and queried, but also interpreted and acted upon by a rich population of logic-based web agents.

\section{For the AI community?}\label{sec:ai-community}

Chapter~7 painted a broad picture of Prolog AI agents on the Agentic Web: symbolic reasoning engines encapsulated as actors, distributed across the network, orchestrating both other symbolic services and neural components such as large language models. In this section we revisit that picture from the vantage point of the wider AI community. If you work on machine learning, symbolic AI, multi-agent systems, or conversational agents, what does the Prolog Trinity ecosystem offer you?

Figure~\ref{fig:prolog-web-intro-20} provides a visual starting point. The central nodes in the diagram host different Prolog systems, but they are surrounded by symbols that clearly point beyond classical logic programming: on the left, a VR headset suggests immersive user interfaces in which Prolog agents act as the cognitive core behind virtual or augmented experiences; on the right, a social robot indicates embodied agents capable of speech, gaze, and gesture, driven by logic-based controllers deployed on the Prolog Web. Near the top, the Python logo together with the OpenAI mark represents a Python-based service that fronts a large language model such as ChatGPT and exposes it as a web API. Web Prolog actors can send prompts to such a service, receive responses, and integrate them with other symbolic knowledge. The same infrastructure can connect to other machine-learning services, data stores, and external tools. The message of the figure is that Prolog agents are meant to inhabit a mixed ecosystem in which traditional AI components, neural models, and interactive front-ends all communicate through a common web-native fabric.

A first benefit for the AI community is that the Prolog Web offers a web-native arena for AI systems. Modern AI already lives on the Web: models are trained on web-scale data, deployed behind web APIs, and consumed through browsers, mobile apps, and embedded devices. The Prolog Web takes this seriously and proposes a programming model in which AI components are not isolated binaries hidden behind opaque endpoints, but named, message-receiving actors inhabiting a shared, logically structured space. For AI researchers interested in reproducible experiments, deployable prototypes, or long-lived services, this means that symbolic components, decision procedures, and coordination mechanisms can be deployed as first-class web entities rather than as ad hoc infrastructure.

Second, the Trinity provides a mature substrate for symbolic and explainable AI. Chapter~7 revisited Prolog's classic strengths in knowledge representation, constraint solving, planning, meta-interpretation, and proof construction. Those strengths become particularly valuable in a world where black-box models dominate: Web Prolog agents can maintain explicit knowledge bases, construct and expose proof trees, encode norms and policies, and mediate between human-understandable rules and data-driven predictions. For AI practitioners facing regulatory or ethical demands for transparency, this makes Web Prolog a serious candidate for the reasoning core of systems that must explain themselves, justify their actions, or enforce formal constraints.

Third, the Prolog Web supports a principled approach to agent-based AI. Toplevels and other actors already form a simple multi-agent substrate; on top of that, more sophisticated agent architectures (BDI-style agents, SCXML statechart actors, conversational web agents) can be built in a language that directly supports logical inference and structured communication. User-interface agents -- embodied conversational agents, dialogue systems, assistants -- can be implemented as Web Prolog actors that combine DCG-based natural language processing with actor-style interaction. Multi-agent-system (MAS) agents can negotiate, allocate resources, and coordinate plans across multiple nodes using the same message-passing primitives. The result is an agent framework in which both UI-agents and MAS-agents can be expressed using a common logical substrate.

At the same time, much remains open at the level of agent abstractions. Chapter~7 also surveyed a number of existing Prolog-derived agent frameworks, such as DALI and LPS, and argued that Web Prolog is a natural host for similar ideas on the Agentic Web. What is still missing in the Trinity ecosystem is a small, well-defined agent framework that is specific to Web Prolog and lives on top of the existing actor and node architecture. Such a framework would have to identify a modest set of concepts -- events, goals, plans, beliefs, norms, perhaps a restricted form of mental state -- and map them cleanly onto ACTOR nodes and their mailboxes, without undermining the simplicity of the underlying concurrency model. It should also take seriously the lessons from previous agent languages and from current initiatives such as the W3C Autonomous Agents on the Web Community Group, which promotes hypermedia MAS aligned with the Web Architecture. From an AI perspective, designing, implementing, and evaluating such a Web-Prolog-native agent layer is an important piece of future work, because it would provide a more user-friendly programming model for building Prolog agents on the Agentic Web and a concrete artefact around which the Prolog, Semantic Web, and AI communities could coordinate their efforts.

Fourth, the Trinity suggests a concrete path toward neuro-symbolic systems. Rather than trying to implement deep learning inside Prolog, the proposal is to treat neural components -- large language models, classifiers, embedding services, vision models -- as callable services that can serve as the definitions of selected predicates. Web Prolog programs then orchestrate these services, enforce symbolic constraints, maintain state, and integrate results with other sources of knowledge. This division of labour respects the strengths of each paradigm: neural models specialise in pattern recognition and language modelling; logic programs specialise in structure, constraint, and explanation. Because LLMs and other neural models are typically exposed via web APIs, they fit naturally into the Prolog Web's communication model, as illustrated by the Python and OpenAI logos in Figure~\ref{fig:prolog-web-intro-20}.

There is also a historical lesson. Expert systems struggled in part because there was no natural infrastructure on which to deploy them, and because knowledge acquisition was labour-intensive. The Web changes both conditions. A web-based agent infrastructure gives symbolic systems a natural deployment platform, and LLM-based tools can ease some aspects of knowledge acquisition, for example by helping to draft rules or transform natural-language descriptions into formal representations that can then be curated and constrained by human experts. In such architectures, Web Prolog can act as the backbone that connects symbolic rules, machine-learned models, and human input into a coherent agentic whole.

Interestingly, there is a very active W3C Community Group dedicated to development of theories and practices for the advancement of more sophisticated web agents.\footnote{\url{https://www.w3.org/community/webagents/}} Here is how they describe their interests and goals:

\begin{quotation}
\textbf{AUTONOMOUS AGENTS ON THE WEB COMMUNITY GROUP}\\
This community group is interested in the design of Web-based Multi-Agent Systems (MAS) for the deployment of world-wide hybrid communities of people and artificial agents on the Web. Our aim is to design a new class of Web-based MAS that are aligned with the Web Architecture to inherit the properties of the Web (world-wide, open, long-lived, etc.), and are also transparent and accountable to support acceptance by people. We are especially interested in the use of Linked Data and Semantic Web standards for weaving a hypermedia fabric that enables uniform interaction among heterogeneous entities: people, artificial agents, devices, digital services, knowledge repositories, etc. We refer to this new class of Web-based MAS as Hypermedia MAS (hMAS). This community group brings together experts actively contributing to advances in autonomous agents and MAS, the Web Architecture and the Web of Things, Semantic Web and Linked Data, and Web standards in general -- as well as any other areas that could contribute to this approach for distributed intelligence on the Web.
\end{quotation}

\noindent It seems to us that the Prolog Trinity ecosystem might be an ideal foundation for building communities of such agents. Its actors are explicitly aligned with the Web Architecture, communicate through standard web protocols, and can be equipped with autonomy in the sense required by Hypermedia MAS: they can follow links, interpret representations, maintain internal state, coordinate with other agents, and respect declaratively specified norms and policies. For the AI community, this means that the abstractions proposed by the W3C group are not merely aspirational. With Web Prolog as the agent language, the Prolog Web as the substrate, and Prolog agents as the computational entities, it becomes possible to experiment with web-scale autonomous agents on a platform that is both logically grounded and architecturally compatible with the Web itself.

For researchers and practitioners in AI, then, what is in it is the opportunity to develop and deploy symbolic, agent-based, and neuro-symbolic systems on a principled web-native foundation. Web Prolog will not replace Python, TensorFlow, or PyTorch, but it can stand alongside them as the logical control layer that ties components together, reasons about their behaviour, and exposes their actions in a form that both humans and other agents can understand.

---

Furhat is a tabletop social robot head developed by the Swedish company Furhat Robotics, designed as a highly expressive platform for face-to-face spoken interaction.\footnote{I am not affiliated with Furhat Robotics, but I know some members of the team personally.} The robot combines back-projected facial animation onto a physical mask with motorised head movements, microphones, speakers and cameras, yielding an embodied conversational agent whose identity (face, voice, personality) can be reconfigured in software. Furhat's software stack includes both a local ``skill'' framework and a realtime WebSocket (WS) API that exposes the robot as a networked resource: external programs written in any langue can connect over WS to send utterances and control non-verbal behaviour, and to receive events such as user speech, gaze and presence. Over this realtime interface it would be straightforward to control every aspect of the conversational behaviour of a Furhat robot from an actor running on an ACTOR node, or from a statechart actor coordinating complex multimodal behaviour, thereby treating the robot as a remote embodied agent within the Prolog Web. As a result, Furhat has become a versatile research platform in human-robot interaction, social robotics and embodied conversational agents, supporting studies in areas such as turn-taking, multi-party interaction, social signalling, education, and healthcare scenarios. The company offers a freely downloadable SDK with a ``virtual Furhat'' for development and research prototyping, while physical robots and advanced cloud services are provided under commercial licences and subscription packages targeted at universities, labs and enterprise deployments.

---

The W3C Machine Learning Working Group [1] focuses on the Web Neural Network API (WebNN) [2] along with an informative report on ethical principles for Web machine learning [3].  

WebNN provides a platform independent API for executing artificial neural network models. More specifically, inference not training, despite the name of the Working Group. The API is currently a W3C Candidate Recommendation. You can try it out on Chrome browsers, but will first need to enable this with a browser flag.

Despite the name of the Working Group, WebNN is actually targeted at inference, not training.  Unlike PyTorch, it omits support for back propagation and stochastic gradient descent.  The existing libraries (Tensorflow and ONNX) are huge, and moreover, lack a high level representation for defining neural network models.


%\section{For the AI community?}
%
%Chapter~7 painted a broad picture of Prolog AI agents on the Agentic Web: symbolic reasoning engines encapsulated as actors, distributed across the network, orchestrating both other symbolic services and neural components such as large language models. In this section we revisit that picture from the vantage point of the wider AI community. If you work on machine learning, symbolic AI, multi-agent systems, or conversational agents, what does the Prolog Trinity ecosystem offer you?
%
%Figure~14.1 provides a visual starting point. The central nodes in the diagram host different Prolog systems, but they are surrounded by symbols that clearly point beyond classical logic programming: on the left, a VR headset suggests immersive user interfaces in which Prolog agents act as the cognitive core behind virtual or augmented experiences; on the right, a social robot indicates embodied agents capable of speech, gaze, and gesture, driven by logic-based controllers deployed on the Prolog Web. Near the top, the Python logo together with the OpenAI mark represents a Python-based service that fronts a large language model such as ChatGPT and exposes it as a web API. Web Prolog actors can send prompts to such a service, receive responses, and integrate them with other symbolic knowledge. The same infrastructure can connect to other machine-learning services, data stores, and external tools. The message of the figure is that Prolog agents are meant to inhabit a mixed ecosystem in which traditional AI components, neural models, and interactive front-ends all communicate through a common web-native fabric.
%
%A first benefit for the AI community is that the Prolog Web offers a web-native arena for AI systems. Modern AI already lives on the Web: models are trained on web-scale data, deployed behind web APIs, and consumed through browsers, mobile apps, and embedded devices. The Prolog Web takes this seriously and proposes a programming model in which AI components are not isolated binaries hidden behind opaque endpoints, but named, message-receiving actors inhabiting a shared, logically structured space. For AI researchers interested in reproducible experiments, deployable prototypes, or long-lived services, this means that symbolic components, decision procedures, and coordination mechanisms can be deployed as first-class web entities rather than as ad hoc infrastructure.
%
%Second, the Trinity provides a mature substrate for symbolic and explainable AI. Chapter~7 revisited Prolog's classic strengths in knowledge representation, constraint solving, planning, meta-interpretation, and proof construction. Those strengths become particularly valuable in a world where black-box models dominate: Web Prolog agents can maintain explicit knowledge bases, construct and expose proof trees, encode norms and policies, and mediate between human-understandable rules and data-driven predictions. For AI practitioners facing regulatory or ethical demands for transparency, this makes Web Prolog a serious candidate for the reasoning core of systems that must explain themselves, justify their actions, or enforce formal constraints.
%
%Third, the Prolog Web supports a principled approach to agent-based AI. Toplevels and other actors already form a simple multi-agent substrate; on top of that, more sophisticated agent architectures (BDI-style agents, SCXML statechart actors, conversational web agents) can be built in a language that directly supports logical inference and structured communication. User-interface agents -- embodied conversational agents, dialogue systems, assistants -- can be implemented as Web Prolog actors that combine DCG-based natural language processing with actor-style interaction. Multi-agent-system (MAS) agents can negotiate, allocate resources, and coordinate plans across multiple nodes using the same message-passing primitives. The result is an agent framework in which both UI-agents and MAS-agents can be expressed using a common logical substrate.
%
%Fourth, the Trinity suggests a concrete path toward neuro-symbolic systems. Rather than trying to implement deep learning inside Prolog, the proposal is to treat neural components -- large language models, classifiers, embedding services, vision models -- as callable services that can serve as the definitions of selected predicates. Web Prolog programs then orchestrate these services, enforce symbolic constraints, maintain state, and integrate results with other sources of knowledge. This division of labour respects the strengths of each paradigm: neural models specialise in pattern recognition and language modelling; logic programs specialise in structure, constraint, and explanation. Because LLMs and other neural models are typically exposed via web APIs, they fit naturally into the Prolog Web's communication model, as illustrated by the Python and OpenAI logos in Figure~14.1.
%
%There is also a historical lesson. Expert systems struggled in part because there was no natural infrastructure on which to deploy them, and because knowledge acquisition was labour-intensive. The Web changes both conditions. A web-based agent infrastructure gives symbolic systems a natural deployment platform, and LLM-based tools can ease some aspects of knowledge acquisition, for example by helping to draft rules or transform natural-language descriptions into formal representations that can then be curated and constrained by human experts. In such architectures, Web Prolog can act as the backbone that connects symbolic rules, machine-learned models, and human input into a coherent agentic whole.
%
%Interestingly, there is a very active W3C Community Group dedicated to development of theories and practices for the advancement of more sophisticated web agents.\footnote{\url{https://www.w3.org/community/webagents/}} Here is how they describe their interests and goals:
%
%\begin{quotation}
%\textbf{AUTONOMOUS AGENTS ON THE WEB COMMUNITY GROUP}\\
%This community group is interested in the design of Web-based Multi-Agent Systems (MAS) for the deployment of world-wide hybrid communities of people and artificial agents on the Web. Our aim is to design a new class of Web-based MAS that are aligned with the Web Architecture to inherit the properties of the Web (world-wide, open, long-lived, etc.), and are also transparent and accountable to support acceptance by people. We are especially interested in the use of Linked Data and Semantic Web standards for weaving a hypermedia fabric that enables uniform interaction among heterogeneous entities: people, artificial agents, devices, digital services, knowledge repositories, etc. We refer to this new class of Web-based MAS as Hypermedia MAS (hMAS). This community group brings together experts actively contributing to advances in autonomous agents and MAS, the Web Architecture and the Web of Things, Semantic Web and Linked Data, and Web standards in general -- as well as any other areas that could contribute to this approach for distributed intelligence on the Web.
%\end{quotation}
%
%It seems to us that the Prolog Trinity ecosystem might be an ideal foundation for building communities of such agents. Its actors are explicitly aligned with the Web Architecture, communicate through standard web protocols, and can be equipped with autonomy in the sense required by Hypermedia MAS: they can follow links, interpret representations, maintain internal state, coordinate with other agents, and respect declaratively specified norms and policies. For the AI community, this means that the abstractions proposed by the W3C group are not merely aspirational. With Web Prolog as the agent language, the Prolog Web as the substrate, and Prolog agents as the computational entities, it becomes possible to experiment with web-scale autonomous agents on a platform that is both logically grounded and architecturally compatible with the Web itself.
%
%For researchers and practitioners in AI, then, what is in it is the opportunity to develop and deploy symbolic, agent-based, and neuro-symbolic systems on a principled web-native foundation. Web Prolog will not replace Python, TensorFlow, or PyTorch, but it can stand alongside them as the logical control layer that ties components together, reasons about their behaviour, and exposes their actions in a form that both humans and other agents can understand.
%
%Chapter~\ref{sec:the-prolog-web-of-ai} argued that Web Prolog is a natural host language for logic-based AI agents and surveyed a number of existing Prolog-derived agent frameworks, such as AgentSpeak, DALI, and related work on logic-based BDI agents. What is still missing in the Trinity ecosystem is a small, well-defined agent framework that is specific to Web Prolog and lives on top of the existing actor and node architecture. Such a framework would have to identify a modest set of concepts -- events, goals, plans, beliefs, norms, and perhaps a restricted form of mental state -- and map them cleanly onto ACTOR nodes and their mailboxes, without undermining the simplicity of the underlying concurrency model. It should also take seriously the lessons from previous agent languages and from current initiatives such as the W3C Autonomous Agents on the Web Community Group, which promotes hypermedia MAS aligned with the Web Architecture. Designing, implementing, and evaluating such a Web-Prolog-native agent layer is therefore an important piece of future work: it would provide a more user-friendly programming model for building Prolog agents on the Agentic Web, and at the same time offer a concrete artefact around which the Prolog, Semantic Web, and AI communities could coordinate their efforts.


\section{For web programmers and the future of web development?}\label{sec:for-web-programmers}

The vast and diverse community of web programmers -- working daily with JavaScript and a wide range of server-side languages and frameworks -- forms much of the backbone of today's digital world. For this community, the Prolog Trinity -- and Web Prolog in particular -- offers a genuinely different way of thinking about certain classes of web applications and services. It is not a call to abandon existing stacks, but an invitation to treat logic as a first-class instrument in web development when the problem domain calls for it.

From a web programmer's perspective, one of the most immediate attractions of Web Prolog is its suitability for rule-intensive applications. Contemporary web stacks are remarkably versatile, yet they can become cumbersome when handling intricate business rules, layered validation logic, sophisticated configuration schemes, or highly personalised behaviour based on many interacting conditions. In such cases, large bodies of imperative code are often needed to simulate what is, at heart, a set of logical constraints and decision rules. Web Prolog offers a more direct route: developers can encode this structure as logical relations and rules, and rely on the inference engine -- with unification and backtracking -- to explore the relevant cases. For certain classes of problems, the resulting code can be more concise, more transparent, and easier to adapt as requirements evolve.

The Trinity also encourages a shift in how we think about the dynamics of web applications. Much of the traditional Web still operates in a client-driven request-response style: a browser or mobile app asks, a server answers. In contrast, the Trinity's long-lived, communicative agents are designed to live continuously on the network. They can subscribe to feeds, monitor data streams, react to events, initiate contact with other services, and cooperate to achieve goals. In the terminology of this book, they inhabit an \emph{Agentic Web}: a web of durable participants and ongoing interactions rather than a web of one-shot calls. For web programmers, this suggests a move towards more proactive, event-driven, and intelligent entities, where some application behaviour is delegated to long-lived agents instead of being recomputed from scratch on each request.

Data integration is another area where Web Prolog can play a useful role. Modern web applications routinely draw on many APIs and heterogeneous data sources, creating ad hoc ``mashups'' by stitching together HTTP requests, JSON transformations, and local caches. Web Prolog, with its relational core and strong support for querying and transforming structured data, offers a natural environment for building more principled integration services. Information from multiple origins can be represented in a uniform logical form, combined using rules, and subjected to explicit constraints. This is also where alignment with RDF-style representations becomes practically useful. The aim is not to replace existing APIs, but to provide a logic-based layer in which their results can be related, checked, and reasoned about.

None of this removes the fact that Prolog has a learning curve, especially for programmers who have spent their careers in imperative or object-oriented paradigms. Unification and backtracking are conceptually different from variables and loops, and it would be naive to pretend otherwise. The Prolog Trinity therefore tries to make the transition as concrete and rewarding as possible. The focus is on web-centric use cases where the benefits of logic programming become visible quickly; on clear, documented APIs that hide some of the low-level details; and on educational resources and playgrounds aimed specifically at web developers, as discussed in earlier chapters. The claim is that, for the right kinds of tasks, the time invested in learning Web Prolog can be repaid in clarity, robustness, and speed of development.

A related theme is the need to make Prolog legible again to the modern developer. The point is not to disguise its identity, but to move beyond the dated picture of Prolog as a niche language from the expert-systems era. The Trinity presents Prolog instead as a language well suited to contemporary concerns: concurrency, distributed agent architectures, web-native communication, and logical intelligence at the edge of applications. For web programmers, the message is that Prolog is no longer confined to standalone AI prototypes or academic examples; it can participate directly in the infrastructure of real web systems.


%It is also worth noting that the broader web community has long been uneasy about relying on a single dominant language. Many developers regard JavaScript as too central. In 2014, Gilad Bracha, then a senior software engineer at Google, argued that web applications might eventually surpass desktop applications in function and usability, but only if developers were given a wider choice of languages.\footnote{\url{https://www.pcworld.com/article/2362500/google-engineer-we-need-more-web-programming-languages.html}} Since then, languages such as Elm and ClojureScript have brought functional programming closer to the browser, expanding the range of styles that can be expressed on the client side. Yet one programming mode remains conspicuously underrepresented in mainstream web development: web-native logic programming. Web Prolog is proposed as one candidate to fill that gap -- not as a replacement for JavaScript or Python, but as a specialised instrument for the logical, rule-driven, and agentive parts of next-generation web applications.

\subsection*{Constructing user interfaces with statecharts}\label{sec:ui-statecharts}

One area in which the Trinity's ideas connect directly to a recognised pain point in web development is user-interface state management. Web applications must ensure that users can perform only the operations that are valid at a given moment, that asynchronous events do not leave the interface in an inconsistent state, and that edge cases -- a button pressed twice, a network timeout during a form submission, a navigation event during a pending request -- are handled gracefully rather than ignored. These are control problems, and they grow harder as applications become more interactive.

The JavaScript ecosystem has responded with a succession of state-management libraries -- Redux, MobX, Vuex, Zustand, and others -- each offering a different trade-off between explicitness and convenience. These tools have improved the situation considerably, but they address the problem at the library level rather than at the level of the underlying model. The developer still reasons about state transitions implicitly, as scattered conditional logic embedded in event handlers and reducers.

Statecharts offer a more structured alternative. As developed in Chapter~\ref{sec:web-prolog-statechart-behavior}, a statechart makes control explicit: the set of possible states, the events that trigger transitions, the guards that constrain them, and the actions that accompany them are all declared as inspectable structure rather than buried in procedural code. This is not a new idea in the web world. In his book \emph{Constructing the User Interface with Statecharts}, Horrocks~(\citeyear{horrocks1999constructing}) argued for statecharts as a design tool for GUIs. More recently, David Khourshid's XState library\footnote{\url{https://github.com/statelyai/xstate}} has brought SCXML-compatible statecharts to the JavaScript front-end community, and his conference talks on the subject have resonated with developers frustrated by ad hoc state management.\footnote{See also the resources collected at \url{https://statecharts.github.io/resources.html}.} David Junger~(\citeyear{jungerclient}) was an early proponent of applying statecharts to web interfaces.

The core insight behind this movement can be stated concisely. A user interface is a reactive system: given a current state and an incoming event, it must produce a new state and a set of actions. A statechart formalises exactly this relationship, decomposing the state space into a hierarchy of named states with explicit transitions, and thereby making the control structure of the interface visible, testable, and communicable to non-programmers. The advantages -- visual design, exhaustive enumeration of modes, decoupling of the state machine from the rendering code -- are precisely what web developers need when applications grow beyond a handful of interactive elements.

Where the Trinity adds something to this picture is in the combination of statecharts with a logic-based scripting and data-modelling language. In the XState world, executable content is JavaScript: guards are functions, actions are functions, and the data model is whatever JavaScript data structures the developer chooses. In the Trinity's statechart actors, the corresponding role is played by Web Prolog. Guard conditions become Prolog goals that succeed or fail; actions become Prolog predicates that can query a knowledge base, enforce constraints, construct structured messages, or invoke remote services; and the data model inherits Prolog's relational and symbolic strengths. For applications in which the interface must reason -- must check policies, validate configurations, maintain and query a structured model of the domain -- this combination is more natural than embedding such logic in JavaScript callbacks.

The practical question is whether web developers would adopt such an approach. The answer depends in part on the availability of Web Prolog in browser environments, a topic discussed in Section~\ref{sec:browser-profile} and acknowledged in Chapter~\ref{sec:from-vision-to-execution} as one of the main open engineering challenges. But it also depends on whether the web development community continues to recognise state management as a fundamental difficulty rather than a solved problem. The evidence from the steady growth of interest in XState and related tools suggests that it does. If statecharts are indeed the right abstraction for governing complex interactive behaviour, then the question of what scripting language sits inside them becomes practically important -- and Prolog's case for that role, in domains where logic and constraint matter, is strong.

More broadly, the relevance of statecharts to web development reinforces a point made throughout this book: that the distinction between ``front-end'' and ``back-end'' obscures a shared underlying problem. Whether one is managing dialogue turns in a voice interface, protocol states in a distributed agent, or navigation flows in a single-page application, the challenge is the same: to govern reactive behaviour so that it remains correct, inspectable, and maintainable as complexity grows. The statechart actor, as developed in Chapter~\ref{sec:web-prolog-statechart-behavior}, is the Trinity's answer to that challenge, and the web front-end is one of the most natural arenas in which to test it.

\section{For the Erlang community?}\label{sec:erlang-community}

Figure~\ref{fig:prolog-web-intro-20} does not display an Erlang logo, yet Erlang is present throughout the entire diagram in a less visible but more fundamental way: without Erlang's intellectual contribution, none of the arrows in the figure would mean what they do. The very possibility of Web Prolog agents communicating asynchronously, isolating their internal state, and coordinating distributed computations derives directly from ideas developed in the Erlang/OTP world. Readers from the Erlang community will recognise the design lineage immediately. What, then, might the Prolog Trinity ecosystem offer in return?

A first answer is that Web Prolog represents a re-engagement with a line of thought that once connected the two communities closely. Erlang's early prototypes grew out of experiments with Prolog, but in order to meet the stringent requirements of fault-tolerant telecom systems, the language shed much of Prolog's logical machinery -- backtracking, unification-heavy control, meta-programming -- in favour of determinism, reliability, and simplicity. The result was a brilliantly focused platform for soft real-time systems. Web Prolog now attempts the complementary move: to reintroduce rich logical capabilities into an actor-style framework that still respects the principles that made Erlang successful.

For Erlang developers, this offers a practical gain. Many applications built atop OTP involve increasingly complex decision-making logic, rich policy constraints, structured plans, or rule-governed behaviour. Encoding such logic in Erlang's functional idiom is certainly possible, but often verbose or difficult to validate. Web Prolog provides an alternative: the ``brains'' of an Erlang-style actor can be written declaratively, using unification, backtracking search, DCGs, constraint solvers, or rule interpreters. In other words, Erlang processes may remain responsible for everything they excel at -- supervision, I/O management, binary protocol handling, predictable latency -- while Web Prolog agents handle symbolic reasoning, explanation, and sophisticated decision procedures.

A second answer concerns interoperability and hybrid architectures. The Prolog Web aims to be a network of services that speak the same actor-oriented protocol. There is no reason an Erlang node could not sit in this network. Indeed, the BEAM VM remains one of the most robust actor engines ever built, and Erlang systems could participate as \emph{first-class} nodes in a heterogeneous set of communicating services. They could delegate logical tasks to Web Prolog agents while providing high-performance endpoints for streaming, front-end connectivity, or massive connection handling. Conversely, Web Prolog agents could be orchestrated or supervised by Erlang systems familiar with OTP's reliability patterns. A shared communication fabric opens the possibility of hybrid systems that neither community has fully explored.

Third, Web Prolog offers a conceptual bridge. If Erlang brought message passing, isolation, linking, and supervision to the world of functional programming, the Prolog Trinity offers Erlang developers a chance to experiment with logical tools that can simplify agent complexity: constraint propagation, rule languages, proof trees, meta-interpreters, planning, explanation, and more. These features can complement, rather than compete with, Erlang's strengths. In this respect, the Trinity plays a role somewhat analogous to Akka for Scala or Cloud Haskell for Haskell -- projects that sought to borrow the essence of Erlang's concurrency model for languages with richer type systems or evaluation strategies. The difference is that Web Prolog was designed from the outset with actor semantics in mind, and its logical features sit \emph{beside} rather than \emph{on top of} the actor substrate.

Finally, there is a research opportunity. Several commentators have noted that while Erlang originally grew out of logic programming ideas, its present form is only one point in a larger design space. Web Prolog explores another. If one were to imagine a future Web Prolog VM that matched the responsiveness and low-latency scheduling of BEAM while still supporting unification, guarded backtracking, and meta-level execution, this would require new research into virtual machine architecture for logic-based actors. That research would draw naturally on decades of BEAM expertise. Conversely, Erlang developers interested in broadening the conceptual reach of Erlang-style concurrency may find in Web Prolog a living laboratory for ideas that cannot easily be explored inside the BEAM constraints.  

In short, the Erlang community has already gifted the wider world a model of concurrency that is elegant, principled, and battle-tested. The Prolog Trinity ecosystem acknowledges that debt and offers something in return: a way to explore richer logic inside an actor discipline, a platform for hybrid systems that combine robust process management with deep symbolic reasoning, and a research frontier where Prolog's declarative core and Erlang's architectural clarity meet again as peers. If successful, this could reopen a line of communication between two communities that once shared foundational ideas -- and may once again have much to offer each other.

